% Ported from thesis/07-threats.md — wording unchanged.
\chapter{Threats to validity \& scope}\label{ch:threats}

% Ported from hypotheses.md "Scope & threats" + scope-of-claims.md.
% Keep the two tables in sync with those sources.

This chapter states the boundary of every claim. We organize the threats in
the standard vocabulary. \textbf{Construct validity} asks whether the measured
quantities mean what the claims need them to mean: whether the probe-based
aggregate score measures ``resilience'' at all is itself under test (H1
quantifies its reliability rather than assuming it), the conntrack flush
percentage is a \emph{proxy} for reconvergence whose causal decomposition is
explicitly pending (\S\ref{sec:h2-mechanism}), and the \texttt{metricAvailability} bookkeeping exists
because a missing metric silently read as zero would corrupt the construct
it stands for. \textbf{Internal validity} asks whether observed differences are
attributable to the manipulated variable: the threats here are run-level
drift, placement mismatch, and dirty provenance, and the defenses --- the
dependent-vs-control route split, within-run correlation,
\texttt{placementMatchRates}, and the \texttt{doctor --strict} gate --- are built into the
design (\S\ref{sec:three-layer}, \S\ref{sec:campaign}) rather than applied post hoc. \textbf{External validity} asks
how far the results carry: a single small virtualized cluster, one
workload, one Kubernetes version and kube-proxy mode, and a single-replica
deployment bound the claims tightly, which is why \S\ref{sec:generalizes} separates what
generalizes (the method, the mechanism's direction) from what does not (the
absolute values, anything multi-replica). \textbf{Conclusion validity} asks
whether the statistics support the inferences drawn: with 7 sessions and
$n = 3$ iterations per cell, the campaign leans on nonparametric tests,
cluster-bootstrap intervals, equivalence testing for the decoupling claims,
and explicit power/MDE calculations rather than significance hunting
(\S\ref{sec:statistics}). The tables below give the threat-by-threat detail.

\section{Threats and defences}\label{sec:threats-table}

\begin{center}
\footnotesize
\begin{tabular}{p{0.2\textwidth}p{0.32\textwidth}p{0.38\textwidth}}
\toprule
\textbf{Threat} & \textbf{Why it matters} & \textbf{Defence} \\
\midrule
\textbf{Single-replica design} & 100\% \texttt{pod-delete} guarantees the only instance disappears, which can swamp topology effects. & Scope the claim to \emph{single-replica churn} and layered measurement; multi-replica anti-affinity is named as future work, not claimed. \\
\addlinespace
\textbf{Small virtualized cluster} & Four 4 GiB KVM/QEMU workers may not generalize. & Claim bounded external validity; report \emph{direction} and \emph{mechanism}, not absolute latency values. \\
\addlinespace
\textbf{Version sensitivity} & kube-proxy / conntrack behaviour evolves across releases (materially in v1.31--v1.32). & Archive exact Kubernetes (v1.28.6), CNI, runtime, kube-proxy mode/conntrack settings, ChaosProbe commit (\texttt{runMetadata} / manifests); present as a measurement study of a specific environment. \\
\addlinespace
\textbf{Placement mismatch} & The scheduler may not realize the intended placement. & Report \texttt{placementMatchRates}; flag or exclude mismatched iterations. \\
\addlinespace
\textbf{Run-to-run drift} & Iteration noise can dominate (H1: 37.6\% run-to-run + 59.1\% iteration variance). & Independent single-commit sessions as unit of analysis; strategy order fixed; pre/post snapshots; run modelled as a random/blocking effect; cluster-bootstrap CIs. \\
\addlinespace
\textbf{Dirty provenance} & Untracked files / missing metadata undermine credibility (the H4 pilot lesson). & Never quote results from runs failing \texttt{doctor --strict}; the dirty H4 pilot was replaced by two doctor-gated $i = 4$ batches. \\
\addlinespace
\textbf{Metric-availability gaps} & Missing PromQL queries can manufacture fake zeros. & \texttt{metricAvailability} distinguishes ``not collected'' from ``collected zero'' (e.g.\ kube-proxy programming latency is recorded as \emph{uncollected} here, and anchors the mechanism only conceptually). \\
\addlinespace
\textbf{Overclaiming causality} & Run-level slowness can confound correlations. & Dependent-vs-control routes + within-run correlation; TOST for equivalence claims; causal language reserved for the manipulated variable (placement). \\
\bottomrule
\end{tabular}
\end{center}

Two of these threats deserve a sentence beyond the table. The
single-replica design is not only a threat but a \emph{scoping instrument}: it
is what makes \texttt{pod-delete} a pure churn fault (the outage is total by
construction, so any placement signal must appear at a non-availability
layer), and it is simultaneously what excludes the multi-replica
anti-affinity regime from every claim. And the run-to-run drift row is not
an admission that the environment was too noisy --- quantifying that noise
and what it does to score-based ranking \emph{is} H1; the campaign design exists
because the drift was measured, not despite it.

\section{What generalizes vs.\ what does not}\label{sec:generalizes}

\begin{center}
\footnotesize
\begin{tabular}{p{0.32\textwidth}p{0.18\textwidth}p{0.4\textwidth}}
\toprule
\textbf{Aspect} & \textbf{Generalizes?} & \textbf{Why / caveat} \\
\midrule
The \emph{method} (placement-aware, cross-layer, provenance-gated chaos evaluation) & \textbf{Yes} & The framework and analysis discipline are reusable on any cluster/workload. \\
\addlinespace
The \emph{direction} of the H2 mechanism effect (spread flushes more conntrack than colocate under churn) & \textbf{Direction only} & Reproducible here; absolute values are environment-specific. \\
\addlinespace
Absolute metric values (latency, recovery s, flush \%) & \textbf{No} & Tied to a small virtualized 5-node cluster. \\
\addlinespace
Mechanism behaviour (conntrack, kube-proxy sync) & \textbf{Environment-contingent} & Depends on CNI, kube-proxy mode (ipvs here), kernel/conntrack settings --- archived per run so the scope is explicit. \\
\addlinespace
The score-instability finding & \textbf{This regime} & Established for single-replica \texttt{pod-delete}; ``cannot rank placement strategies under session variance'', not ``scores don't work''; not asserted for multi-replica or contention-dominated regimes. \\
\addlinespace
The H5 separator & \textbf{As a validated static screen here} & Two-regime separation replicated across two independent batches ($\rho$ 0.79 $\rightarrow$ 0.25: the continuous correlation did not replicate and is not claimed); validation of the metric, not of the locality concept. \\
\addlinespace
The H6 trade-off & \textbf{Qualitatively known; quantified here} & Two-point contrast, single-replica; the \emph{quantification} is this environment's. \\
\addlinespace
Anything about multi-replica / HA failover & \textbf{Not in scope} & Structurally excluded by the single-replica design (\S\ref{sec:future-work}). \\
\bottomrule
\end{tabular}
\end{center}

The asymmetry of this table is deliberate. The portable row is the first
one: the measurement design, the campaign protocol, and the analysis
scripts transfer to any cluster and workload unchanged, and a replication
on different infrastructure is the most valuable follow-up this study
could receive. Everything below it degrades gracefully from ``direction
only'' to ``not in scope'' --- and where a row says environment-contingent, the
archived manifests make the contingency checkable rather than vague: a
reader who wants to know whether a result could apply to their cluster can
compare their kube-proxy mode, conntrack settings, and Kubernetes version
against the archived fingerprint instead of guessing.

\section{Explicitly not claimed}\label{sec:not-claimed}

Per \texttt{scope-of-claims.md}:
no universal ``best'' placement strategy; no ``spread is worse/disproven''; no
proven causality for the conntrack mechanism (mechanistic consistency, not
controlled causal proof); no generalization to Kubernetes clusters broadly;
no ``the score identifies the best strategy''; no unqualified ``reproducible''
(only the mechanism metrics reproduce, and only with the archived rerun
package).

These exclusions are commitments, not disclaimers: each one names a stronger
claim that the data could be mistaken to support and that we decline. The
list was not assembled at writing time; it accumulated during the study, as
earlier, stronger phrasings (``refuted'', ``spread is disproven'', a $\sim$3$\times$ user-layer
advantage) failed replication or scrutiny and were retired (\S\ref{sec:h4},
\S\ref{sec:l1l3}). Where this document's prose and that scope statement could ever be
read to disagree, the scope statement governs.
